\subsection{ms-cursos --- Cursos, secciones, matrícula y anuncios}
\label{sec:ms-cursos}
\textit{Responsable: Benjamín y Edinson (ACA-03, ACA-04, COL-01). Repositorio: \url{https://github.com/nubeduca/nubeduca-ms-cursos}.}

ms-cursos gestiona los cursos de cada periodo, sus secciones, quién está matriculado y los anuncios. Un curso nace en borrador (\texttt{DRAFT}), se publica (\texttt{PUBLISHED}) cuando los alumnos ya pueden verlo y se archiva (\texttt{ARCHIVED}) para dejarlo de solo lectura; cada transición emite su evento y sube la \texttt{aggregateVersion} del curso (1, 2 y 3), así los consumidores descartan los eventos que llegan en desorden. Solo se crean cursos en un periodo activo, y el código del curso (por ejemplo \texttt{CS2032}) es único dentro del periodo gracias a un ítem de reserva \texttt{COURSECODE\#} que va en la misma transacción.

El servicio no llama a ms-usuarios: mantiene dos proyecciones locales alimentadas por eventos. Con \texttt{term.started} guarda qué periodos están activos, y con \texttt{user.provisioned} guarda el rol y el código institucional de cada usuario. Con esa segunda proyección funciona la matrícula por CSV (\texttt{institutionalCode,role,section}, hasta 2\,000 filas): cada fila válida crea una matrícula y emite un \texttt{enrollment.created}, y las filas con error se devuelven con su número de línea y el motivo, sin detener la importación. Cada matrícula se guarda tres veces en una sola transacción: por curso (\texttt{ENROLL\#}, para listar los matriculados), por id (\texttt{ENROLLID\#}, para el \texttt{DELETE}) y por usuario (\texttt{USERCOURSE\#}, para los cursos del alumno), de modo que cada consulta es una \texttt{query} por prefijo. Los anuncios se ordenan por su ULID, del más reciente al más antiguo, y su evento lleva solo el título: el cuerpo se lee con \texttt{GET}.

El tenant y el rol salen siempre del token. El administrador ve y gestiona todos los cursos de su tenant. El docente crea cursos y solo puede editar, publicar, archivar, crear secciones y matricular en los suyos; para publicar un anuncio, además, tiene que estar matriculado como \texttt{TEACHER} en el curso, y el curso debe estar publicado. El alumno solo ve los cursos publicados donde está matriculado, con sus secciones y anuncios; a un curso sin acceso se le responde 404, para no revelar que existe.

\subsubsection{Endpoints}
\begin{table}[htbp]
\centering\footnotesize
\begin{tabularx}{\textwidth}{lXll}
\toprule
Método & Ruta & Roles & Emite \\
\midrule
GET    & \texttt{/health} & pública & --- \\
GET    & \texttt{/courses} & todos & --- \\
POST   & \texttt{/courses} & TEACHER, ADMIN & \texttt{course.created} \\
GET    & \texttt{/courses/\{courseId\}} & todos & --- \\
PATCH  & \texttt{/courses/\{courseId\}} & TEACHER, ADMIN & --- \\
POST   & \texttt{/courses/\{courseId\}/publish} & TEACHER, ADMIN & \texttt{course.published} \\
POST   & \texttt{/courses/\{courseId\}/archive} & TEACHER, ADMIN & \texttt{course.archived} \\
GET    & \texttt{/courses/\{courseId\}/sections} & todos & --- \\
POST   & \texttt{/courses/\{courseId\}/sections} & TEACHER, ADMIN & --- \\
GET    & \texttt{/courses/\{courseId\}/enrollments} & TEACHER, ADMIN & --- \\
POST   & \texttt{/courses/\{courseId\}/enrollments} & TEACHER, ADMIN & \texttt{enrollment.created} \\
POST   & \texttt{/courses/\{courseId\}/enrollments/import} & TEACHER, ADMIN & \texttt{enrollment.created} \\
DELETE & \texttt{/courses/\{courseId\}/enrollments/\{enrollmentId\}} & TEACHER, ADMIN & \texttt{enrollment.dropped} \\
GET    & \texttt{/courses/\{courseId\}/announcements} & todos & --- \\
POST   & \texttt{/courses/\{courseId\}/announcements} & TEACHER & \texttt{announcement.published} \\
\bottomrule
\end{tabularx}
\caption{Endpoints de ms-cursos (\texttt{apis/ms-cursos.openapi.yaml} de \texttt{nubeduca-contratos}); \textit{todos} = STUDENT, TEACHER y ADMIN.}
\end{table}

\subsubsection{Eventos}
\textbf{Publica:}
\begin{itemize}
  \item \texttt{lms.courses.course.created.v1}
  \item \texttt{lms.courses.course.published.v1}
  \item \texttt{lms.courses.course.archived.v1}
  \item \texttt{lms.courses.enrollment.created.v1}
  \item \texttt{lms.courses.enrollment.dropped.v1}
  \item \texttt{lms.courses.announcement.published.v1}
\end{itemize}

\textbf{Consume:}
\begin{itemize}
  \item \texttt{lms.users.term.started.v1}
  \item \texttt{lms.users.user.provisioned.v1}
\end{itemize}

\paragraph{Hallazgo de la prueba de carga (PERF-01).}
\texttt{GET /courses} hace \texttt{query} de toda la partición del tenant y filtra \texttt{ownerId} con
\texttt{FilterExpression}, sin índice: con el seed de 30\,000 tarda unos 1{,}7\,s (p95 1{,}84\,s)
incluso con un usuario, mientras que \texttt{GET /courses/\{courseId\}} responde en 43\,ms (p95). Un
GSI por \texttt{ownerId} lo dejaría en decenas de milisegundos (sección~\ref{sec:pruebas}).

El índice no se agregó antes de la entrega: la tabla de \texttt{serverless.yml} en \texttt{main} sigue
sin GSI, así que la latencia de \texttt{GET /courses} es la medida en PERF-01. Agregar el GSI por
\texttt{ownerId} queda como trabajo futuro (sección~\ref{sec:conclusiones}).

\subsubsection{Estructura JSON de la tabla \texttt{<stage>-t\_cursos}}
Clave de partición \texttt{tenant\_id}, clave de ordenamiento \texttt{sk}.

\begin{table}[htbp]
\centering\small
\begin{tabularx}{\textwidth}{>{\raggedright\arraybackslash}p{0.46\textwidth}X}
\toprule
\texttt{sk} & Qué es \\
\midrule
\texttt{COURSE\#<courseId>} & Curso \\
\texttt{COURSECODE\#<termId>\#<code>} & Reserva del código del curso en el periodo \\
\texttt{SECTION\#<courseId>\#<sectionId>} & Sección del curso \\
\texttt{ENROLL\#<courseId>\#<userId>} & Matrícula \\
\texttt{ENROLLID\#<enrollmentId>} & Puntero a la matrícula, para el \texttt{DELETE} \\
\texttt{USERCOURSE\#<userId>\#<courseId>} & Cursos de un usuario \\
\texttt{ANN\#<courseId>\#<announcementId>} & Anuncio \\
\texttt{PROY\#TERM\#<termId>} & Periodo activo (proyección de \texttt{term.started}) \\
\texttt{PROY\#USER\#<userId>} & Rol y código del usuario (proyección de \texttt{user.provisioned}) \\
\texttt{PROY\#USERCODE\#<código>} & Código institucional $\rightarrow$ usuario, para el CSV \\
\texttt{EVENTO\#}, \texttt{INBOX\#}, \texttt{AGG\#} & Outbox, inbox y versión de agregados de la plantilla \\
\bottomrule
\end{tabularx}
\caption{Prefijos de \texttt{sk} en la tabla de ms-cursos.}
\end{table}

\begin{lstlisting}
{
  "tenant_id": "utec",
  "sk": "COURSE#crs_00000000000000000000000020",
  "entidad": "COURSE",
  "version": 1,
  "courseId": "crs_00000000000000000000000020",
  "termId": "trm_0000000000000000000000SEED",
  "code": "SEED-00020",
  "title": "Cloud Computing 3",
  "description": "Curso de seed para pruebas de carga.",
  "status": "PUBLISHED",
  "ownerId": "3651db88-79cd-5df8-bd8c-488941b6e878",
  "createdAt": "2026-08-17T13:00:00Z",
  "seed": true
}
\end{lstlisting}
Ítem real leído de la tabla desplegada \texttt{dev-t\_cursos} con \texttt{aws dynamodb get-item} (22-Nov-2026).

Los demás tipos de ítem muestran los atributos que escribe el código de \texttt{api/} y
\texttt{consumidores/}, con un curso, un docente y la alumna demo Ana Torres de ejemplo. La reserva
del código y el puntero de la matrícula:

\begin{lstlisting}
{
  "tenant_id": "utec",
  "sk": "COURSECODE#trm_01JCP0KWJDXF9GJ9YNTK2H84A9#CS2032",
  "entidad": "COURSECODE",
  "courseId": "crs_01JCQRZHQT9F4WY2JYPDV546K7"
}
\end{lstlisting}

\begin{lstlisting}
{
  "tenant_id": "utec",
  "sk": "ENROLLID#enr_01JCR1B7Y2K8M4Q6T9V3W5X0Z1",
  "entidad": "ENROLLMENT_ID",
  "courseId": "crs_01JCQRZHQT9F4WY2JYPDV546K7",
  "userId": "11111111-1111-4111-8111-111111111111"
}
\end{lstlisting}

Sección, matrícula y su copia por usuario:

\begin{lstlisting}
{
  "tenant_id": "utec",
  "sk": "SECTION#crs_01JCQRZHQT9F4WY2JYPDV546K7#sec_01JCR0A4T7Q9W2N8M5K3H6G1F0",
  "entidad": "SECTION",
  "courseId": "crs_01JCQRZHQT9F4WY2JYPDV546K7",
  "sectionId": "sec_01JCR0A4T7Q9W2N8M5K3H6G1F0",
  "name": "Sección 1",
  "createdBy": "c62d9930-b473-4b1d-adec-72fdda7cf599",
  "createdAt": "2026-11-13T15:10:00Z"
}
\end{lstlisting}

\begin{lstlisting}
{
  "tenant_id": "utec",
  "sk": "ENROLL#crs_01JCQRZHQT9F4WY2JYPDV546K7#11111111-1111-4111-8111-111111111111",
  "entidad": "ENROLLMENT",
  "courseId": "crs_01JCQRZHQT9F4WY2JYPDV546K7",
  "version": 1,
  "enrollmentId": "enr_01JCR1B7Y2K8M4Q6T9V3W5X0Z1",
  "userId": "11111111-1111-4111-8111-111111111111",
  "role": "STUDENT",
  "sectionId": "sec_01JCR0A4T7Q9W2N8M5K3H6G1F0",
  "createdAt": "2026-11-14T16:00:00Z"
}
\end{lstlisting}

\begin{lstlisting}
{
  "tenant_id": "utec",
  "sk": "USERCOURSE#11111111-1111-4111-8111-111111111111#crs_01JCQRZHQT9F4WY2JYPDV546K7",
  "entidad": "USER_COURSE",
  "courseId": "crs_01JCQRZHQT9F4WY2JYPDV546K7",
  "userId": "11111111-1111-4111-8111-111111111111",
  "role": "STUDENT",
  "enrollmentId": "enr_01JCR1B7Y2K8M4Q6T9V3W5X0Z1"
}
\end{lstlisting}

Anuncio:

\begin{lstlisting}
{
  "tenant_id": "utec",
  "sk": "ANN#crs_01JCQRZHQT9F4WY2JYPDV546K7#ann_01JCS2D9F4H6K8N0P2R4T6V8X0",
  "entidad": "ANNOUNCEMENT",
  "courseId": "crs_01JCQRZHQT9F4WY2JYPDV546K7",
  "announcementId": "ann_01JCS2D9F4H6K8N0P2R4T6V8X0",
  "title": "Bienvenidos al curso",
  "body": "El primer laboratorio es el lunes.",
  "authorId": "c62d9930-b473-4b1d-adec-72fdda7cf599",
  "publishedAt": "2026-11-15T14:00:00Z"
}
\end{lstlisting}

Proyecciones de los eventos de ms-usuarios:

\begin{lstlisting}
{
  "tenant_id": "utec",
  "sk": "PROY#TERM#trm_01JCP0KWJDXF9GJ9YNTK2H84A9",
  "entidad": "PROYECCION",
  "termId": "trm_01JCP0KWJDXF9GJ9YNTK2H84A9",
  "name": "2026-2",
  "startsOn": "2026-08-17",
  "endsOn": "2026-12-12",
  "status": "ACTIVE",
  "recibidoDe": "lms.users",
  "version": 2
}
\end{lstlisting}

\begin{lstlisting}
{
  "tenant_id": "utec",
  "sk": "PROY#USER#11111111-1111-4111-8111-111111111111",
  "entidad": "PROYECCION",
  "userId": "11111111-1111-4111-8111-111111111111",
  "role": "STUDENT",
  "institutionalCode": "202310457",
  "recibidoDe": "lms.users",
  "version": 1
}
\end{lstlisting}

\begin{lstlisting}
{
  "tenant_id": "utec",
  "sk": "PROY#USERCODE#202310457",
  "entidad": "PROYECCION",
  "userId": "11111111-1111-4111-8111-111111111111",
  "role": "STUDENT"
}
\end{lstlisting}

\subsubsection{Datos de prueba}
El 20-Nov-2026, en la cuenta de Dev 2 (stage \texttt{dev}), \texttt{poblar\_datos.py}
cargó \textbf{30\,000 cursos} en \texttt{dev-t\_cursos}: UTEC 15\,000, UPC 10\,000 y PUCP 5\,000. Cada tenant
recibe un periodo activo de seed y 20 docentes de seed; los estados quedan en 60\,\% PUBLISHED, 30\,\% DRAFT y
10\,\% ARCHIVED. Un conteo de los ítems \texttt{COURSE\#} de UTEC devolvió 15\,000.

El 22-Nov-2026 se repitió la carga en otra cuenta (stage \texttt{dev}, cuenta de Dev 4):

\begin{lstlisting}
$ python scripts/poblar_datos.py --stage dev
utec: 15000 cursos
upc: 10000 cursos
pucp: 5000 cursos
Total: 30000 cursos en dev-t_cursos
\end{lstlisting}

El conteo de los ítems \texttt{COURSE\#} de los tres tenants, con \texttt{aws dynamodb query ... --select COUNT}, confirma UTEC 15\,000, UPC 10\,000 y PUCP 5\,000 (figura~\ref{fig:ms-cursos-datos}).

\evidencia{ms-cursos-datos.png}{ms-cursos: conteo de cursos por tenant en \texttt{dev-t\_cursos}.}{fig:ms-cursos-datos}

\subsubsection{Resultados y evidencias}
\evidencia{ms-cursos-evidencia.png}{ms-cursos desplegado en AWS: \texttt{GET /courses/\{courseId\}} como administrador de UTEC devuelve 200 con el primer curso del seed.}{fig:ms-cursos}
